% Pursuing Stacks (cotes 134-2 to 134-8): archive cote 134-2, site page 72 (PDF page 73).
% appendix (letters to Breen). Draft compared: ps1.tex:1577-1605.
% AI-transcribed by Claude Opus 5.5 (claude-opus-5-5, adaptive thinking, effort high), Message Batches API, batch ps1 (msgbatch_01Rv5hCaNTtUt4JHx1yqoHPL), 2026-11-12; not yet checked by eye.
\scanpage{134-2/72}
\mspage{- \unsure{57} -}
\[ F_0 \ , \ \widetilde{F}_1 \ , \ \widetilde{u} : p^*(F_0) \to j^*(\widetilde{F}_1) \]
where $F_0(\widetilde{F}_1)$ are locally constant sheaves on $X_0$ (respectively $X_1$).
It is necessary to recall that \deleted{up to} here $p$ is a real fibration, and
$j$ is an inclusion (in practice, \deleted{``}for the case $X_0$ , $X_1$ smooth\hand{,}\deleted{''} the
inclusion of the boundary in a manifold with boundary).

If you prefer, one can also take the diagram which is less pretty (but
a little more canonical)
\diagram{\[\begin{tikzcd}
  \overset{\circ}{T} \ar[r,"j'"]\ar[d,swap,"p'"] & X_1 \\
  X_0 &
\end{tikzcd}\]}
coming essentially to the same thing, as it is formed from spaces homotopic
to the preceding one. One can even replace $X_0$ by $T$ ($X_0$ being a
deformation retract of it) and write
\diagram{\[\begin{tikzcd}
  \overset{\circ}{T} \ar[r]\ar[d,swap,"p''"] & X_1 \\
  T &
\end{tikzcd}\]}
where ``literally'' $p''$ is now an inclusion, but ``morally'', it is a \textit{fibration}
with \deleted{a} very pretty fibres [notably \inserted{compact} of finite dimension, and \inserted{moreover} \deleted{in fact perhaps}
non-singular varieties - this is \deleted{\unsure{then}} much better than that which is given
by the yoga of Cartan-Serre ``every continuous mapping is equivalent to a
fibration''...] . This last diagram \deleted{is} \inserted{however has the advantage of being} amenable to a purely algebraic, direct
construction, in the context of schemes, once one has developed the
construction of étale tubular neighbourhood\hand{s}\hand{\textsuperscript{(28)}}.
\margin{\unsure{à part} (circled), with a hand-drawn line under the paragraph above and arrows pointing up and down in the left margin}
\inserted{(App \unsure{17})} The point I wish to come to, is that the consideration of sheaves
(or complexes thereof, or n-stacks ...) which are P-constructible on an $X$ ,
where $P$ is a given ``equi-singular'' stratification, \deleted{comes} \inserted{reduces} in our particular
cases to the knowledge of a diagram of ordinary \textit{homotopy types} (or pro-types,
if one comes back to the étale topology)

\margin{72}
